\documentclass[a4paper,twoside]{article}

\usepackage{morewrites}

\usepackage[svgnames,dvipsnames,x11names]{xcolor}
\usepackage{graphicx}
\usepackage{texsmith-lists}
\usepackage[english]{babel}
\usepackage[colorlinks=true, linkcolor=NavyBlue, urlcolor=NavyBlue, citecolor=NavyBlue]{hyperref}
\usepackage[a4paper]{geometry}
\usepackage{ts-fonts}
\usepackage{amsmath}
\usepackage{float}
\usepackage{seqsplit}
\usepackage{ts-code}

\usepackage{lastpage}
\usepackage{refcount}
\makeatletter
\def\ps@plain{
\let\@oddhead\@empty
\let\@evenhead\@empty
\def\@oddfoot{
\hfil
\ifnum\getpagerefnumber{LastPage}>1\relax
\thepage
\fi
\hfil}
\let\@evenfoot\@oddfoot
}
\makeatother

\title{Fonts and scripts}
\author{}\date{}
\begin{document}
\maketitle
TeXSmith automatically wraps non‑Latin scripts in moving arguments (headings, captions, index entries, ...) with per‑script font macros, using a cached lookup built from the Noto family. This gives fast, consistent multilingual output without manual fontspec boilerplate.
\section{Inspect detected fallback fonts}
Use the render command with \tscodeinline{-\allowbreak{}-\allowbreak{}fonts-\allowbreak{}info} to display the scripts and fonts that were detected during a build. The flag works for both direct \LaTeX{} output and full template renders:
\begin{tscode}[lang=bash, engine=pygments]
\begin{Verbatim}[commandchars=\\\{\},breaklines, breakanywhere, commandchars=\\\{\}]
uv\PY{+w}{ }run\PY{+w}{ }texsmith\PY{+w}{ }examples/dialects/dialects.md\PY{+w}{ }\PYZhy{}\PYZhy{}template\PY{+w}{ }article\PY{+w}{ }\PYZhy{}\PYZhy{}fonts\PYZhy{}info
\end{Verbatim}
\end{tscode}
The report includes:
\begin{itemize}
\item the script name and the generated \LaTeX{} commands (\tscodeinline{\textbackslash{}text<slug>} / \tscodeinline{\textbackslash{}<slug>font});
\item the font family chosen for that script (from the cached fallback index);
\item the number of codepoints seen for that script in the current render.
\end{itemize}
When Rich is available, the information is shown as a table; otherwise a plaintext list is printed. The option is non-intrusive and does not change the output artefacts.
\end{document}
